(1) 在随环转动的参考系，用单自由度拉氏函数消去已给定的方位运动，\(T=mR^2(\dot\theta^2+\Omega^2\sin^2\theta)/2\)，\(V=mgR(1-\cos\theta)\)。所以 \(U_{\rm eff}=mgR(1-\cos\theta)-m\Omega^2R^2\sin^2\theta/2\)，方程为 \(mR^2\ddot\theta+\partial_\theta U_{\rm eff}=0\)，即 \(\ddot\theta+\sin\theta(g/R-\Omega^2\cos\theta)=0\)。
(2) 驻点满足 \(\sin\theta(g/R-\Omega^2\cos\theta)=0\)。最低点 \(\theta=0\) 和最高点 \(\theta=\pi\) 始终存在；当 \(\Omega^2R>g\) 时又有对称分支 \(\theta_\pm=\pm\arccos[g/(\Omega^2R)]\)。由 \(U_{\rm eff}''\) 的符号，最低点在阈值前稳定而阈值后失稳；最高点一直不稳定，两个偏转分支在阈值后稳定，呈超临界叉形分岔。
(3) 最低点附近 \(U_{\rm eff}''(0)=mR(g-\Omega^2R)\)，从而 \(\omega_0^2=g/R-\Omega^2\)。偏转支上 \(\cos\theta_* =g/(\Omega^2R)\)，曲率给出 \(\omega_*^2=\Omega^2[1-g^2/(\Omega^4R^2)]\)。在临界点二阶项消失，展开势中的正四次项决定缓慢非线性回复；因此不能把临界频率为零误解为有限周期正弦振荡。
(4) 阻尼方程改为 \(mR^2\ddot\theta+\gamma\dot\theta+\partial_\theta U_{\rm eff}=0\)；对时间求能量导数得 \(d[mR^2\dot\theta^2/2+U_{\rm eff}]/dt=-\gamma\dot\theta^2\le0\)，这里 \(\Omega\) 恒定且驱动器供给维持转速所需的功。过临界时对称的两极小同时出现，最终分支由微小扰动而非理想对称方程确定。
